% !PRINTED 571
\weberpage{571}
%<<CONT
and if we thus denote the variables of $\Phi$ by $t, u_1, u_2,
\ldots$, we shall also set
\weq{5}{\Phi(t) = \Phi(t, u_1, u_2, \ldots),}
or more briefly $\Phi(t, u)$, and this function contains integral
rational numerical coefficients. If we raise the function $\Phi(t)$
to the $p$th power and apply the formula \S.\,150, 4., there follows
from the congruence (3) a new congruence
\weq{6}{\Phi(\tau^p, u_1^p, u_2^p, \ldots) \equiv 0\ (\mathrm{mod}\ \mathfrak p).}

But in the same way we also have
\weq{7}{\tau^p \equiv \omega_1^p t_1^p + \omega_2^p t_2^p + \cdots + \omega_n^p t_n^p.}

If we substitute this in the congruence (6), there arises a
function divisible by $\mathfrak p$, in which the variables
$u_1, u_2, \ldots$, among which indeed the $t_1, t_2, \ldots, t_n$
are included, occur only in the $p$th power. The congruence must
therefore remain correct if we replace the $u_1^p, u_2^p,
\ldots$ and hence also the $t_1^p, t_2^p, \ldots, t_n^p$ by
independent variables $u_1, u_2, \ldots t_1, t_2, \ldots, t_n$
(\S.\,143).

If we accordingly set
\weq{8}{\tau_1 = \omega_1^p t_1 + \omega_2^p t_2 + \cdots + \omega_n^p t_n,}
then from (6) there results
\[
\Phi(\tau_1, u_1, u_2, \ldots) \equiv 0\ (\mathrm{mod}\ \mathfrak p),
\]
and consequently $\tau_1$ too is a root of the congruence (4).

This same procedure can be applied repeatedly, and we find
that
\[
\tau_2 = \omega_1^{p^2} t_1 + \omega_2^{p^2} t_2 + \cdots + \omega_p^{p^2} t_n % [sic: Weber prints \omega_p here; (9) has \omega_n.]
\]
is also a root of the congruence (4), and so on.

If, then, we
define a system of forms $\tau_r$ by
\weq{9}{\tau_r = \omega_1^{p^r} t_1 + \omega_2^{p^r} t_2 + \cdots + \omega_n^{p^r} t_n}
for arbitrary positive exponents $r$, all these quantities
$\tau_r$ are at the same time roots of the congruence (4). It
remains to answer the question how many of these forms $\tau_r$
are different from one another.

By \S.\,150, (4) it is certain that
\[
\tau_r \equiv \tau_{r'}\ (\mathrm{mod}\ \mathfrak p),
\]
if
\[
r \equiv r'\ (\mathrm{mod}\ f)
\]
holds, and consequently among the $\tau_r$ there are certainly not
more than $f$ that are different with respect to the modulus $\mathfrak p$
\weq{10}{\tau, \tau_1, \tau_2, \ldots, \tau_{f-1}.}
